%\documentclass[12pt,oneside]{amsart}
\documentclass{scrartcl}

%\documentclass{article}

\usepackage{amsthm,amsmath,amssymb}
\usepackage[numbers]{natbib}

%\usepackage[colorlinks]{hyperref}
\usepackage{hyperref}
\hypersetup{
    colorlinks,%
    citecolor=black,%
    filecolor=black,%
    linkcolor=black,%
    urlcolor=black
}
\usepackage{hypernat}
\usepackage[top=2.5cm, bottom=2.5cm, left=2.8cm,
right=2.8cm]{geometry}
\usepackage{graphicx}


%\setlength{\parskip}{0pt}
%\setlength{\parsep}{0pt}
%\setlength{\headsep}{0pt}
%\setlength{\topskip}{0pt}
%\setlength{\topmargin}{0pt}
%\setlength{\topsep}{0pt}
%\setlength{\partopsep}{0pt}

\linespread{1}

%\usepackage{marvosym}

%\textwidth = 6.5 in \textheight = 9.2 in \oddsidemargin = 0.0 in
%\evensidemargin = 0.0 in \topmargin = -0.0 in \headheight = 0.0 in
%\headsep = 0.0in \parskip = 0.0in \parindent = 0.0in
%\def\baselinestretch{0.871}


\newtheorem{theorem}{Theorem}[section]
\newtheorem{proposition}[theorem]{Proposition}
\newtheorem{problem}[theorem]{Problem}
\newtheorem{fact}[theorem]{Fact}
\newtheorem{lemma}[theorem]{Lemma}
\newtheorem{defn}[theorem]{Definition}
\newtheorem{corollary}[theorem]{Corollary}
\newtheorem{remark}{Remark}[section]
\newtheorem{example}[theorem]{Example}

\newcommand{\E}{{\mathbb E}}
\newcommand{\se}{{\mathcal{E}}}
\newcommand{\R}{{\mathbb R}}
\renewcommand{\P}{{\mathbb P}}
\newcommand{\ac}{{\mathcal{A}}}
\newcommand{\A}{{\mathcal{A}}}
\newcommand{\B}{{\mathcal{B}}}
\newcommand{\C}{{\mathcal{C}}}
\newcommand{\M}{{\mathcal{M}}}
\newcommand{\Mf}{{\mathfrak{M}}}
\newcommand{\T}{{\mathcal{T}}}
\newcommand{\Z}{{\mathcal{Z}}}
\newcommand{\K}{{\mathcal{K}}}
\newcommand{\lc}{{\mathcal{L}}}
\newcommand{\Ps}{{\mathcal{P}}}
\newcommand{\G}{{\mathcal{G}}}
\newcommand{\pa}{{\mathring{p}}}
\newcommand{\F}{{\cal F}}
\newcommand{\samp}{{\mathcal{X}}}
\newcommand{\X}{{\mathcal{X}}}
\newcommand{\hi}{{\hat{\Phi}}}
\newcommand{\data}{{\text{data}}}
\newcommand{\relu}{{\text{ReLU}}}

\newcommand{\ridge}{{\hat{\beta}_{\text{ridge}}(\lambda)}}
\newcommand{\lasso}{{\hat{\beta}_{\text{lasso}}(\lambda)}}
\newcommand{\ridgemu}{{\hat{\mu}_t^{\text{ridge}}(\lambda)}}
\newcommand{\lassomu}{{\hat{\mu}_t^{\text{lasso}}(\lambda)}}


\newcounter{rcnt}[section]
\renewcommand{\thercnt}{(\roman{rcnt})}

\def\qt#1{\qquad\text{#1}}

\def\argmin{\mathop{\rm argmin}}
\def\argmax{\mathop{\rm argmax}}


\setlength{\parskip}{1.4 \medskipamount}
%\sloppy
%\linespread{1.3}

\begin{document}

\title{STAT 238 - Bayesian Statistics  \\
  Lecture Twenty Two} 
\subtitle{Spring 2026, UC Berkeley}

\author{Aditya Guntuboyina}


\date{16 March 2026}

\maketitle

\section{Gaussian Processes}
The goal is to infer an unknown function $f : \Omega \rightarrow \R$
that is defined on a known domain $\Omega \subseteq \R^p$. We assume
that $\{f(x), x \in \Omega\}$ forms a Gaussian process with mean zero
and covariance function or \textit{kernel} given by $K(x, x')$ i.e., 
\begin{equation*}
  \text{Cov}(f(x), f(x')) = K(x, x') ~~ \text{for all $x, x' \in
    \Omega$}. 
\end{equation*}
The kernel needs to be positive semi-definite i.e., for every $N \geq
1$, distinct points $u_1, \dots, u_N \in \Omega$, the $N \times N$
matrix with $(i, j)^{th}$ entry $K(u_i, u_j)$ is positive 
semi-definite. Sometimes, this matrix will be positive definite, and
hence invertible.

Here are some examples of Gaussian processes and kernels.
\begin{enumerate}
\item \textbf{Brownian Motion}: Here $\Omega = [0, \infty)$ and $K(s,
  t) = \min(s, t)$.
\item \textbf{(scaled) Brownian Motion plus constant}: Here again $\Omega =
  [0, \infty)$. In the Brownian motion model, $f(0) = 0$ which is
  generally an unrealistic assumption. To fix this, we assume that
  $f(t) = \beta_0 + \tau W_t$ where $W_t 
  \sim BM$, $\beta_0 \sim N(0, C)$ and $\tau > 0$ (and independence
  between   $\beta_0$ and $\{W_t\}$). $C$ will be taken to be
  large (potentially $C \rightarrow \infty$). Now  
  \begin{equation*}
    K(s, t) = C + \tau^2 \min(s, t). 
  \end{equation*}
\item \textbf{Integrated Brownian Motion}: $\Omega = [0,
  \infty)$ and
  \begin{equation*}
    f(t) = \int_0^t W_s ds ~~ \text{where $W_s \sim BM$}. 
  \end{equation*}
  The kernel is given by
  \begin{align*}
    K(s, t) &= \text{Cov} \left(\int_0^s W_u du, \int_0^t W_v dv
              \right) \\
            &= \int_0^s \int_0^t \text{Cov}(W_u, W_v) dv du 
            = \int_0^s \int_0^t \min(u, v) dv du
  \end{align*}
  To simplify further, assume that $s \leq t$ so that
  \begin{align*}
K(s, t)            &= \int_0^s \left(\int_0^u v dv + \int_u^t u dv \right) du
    \\
            &= \int_0^s \left(\frac{u^2}{2} + u(t - u) \right) du = \int_0^s \left(ut - \frac{u^2}{2} \right) du = t \frac{s^2}{2}
      - \frac{s^3}{6}. 
  \end{align*}
  For general $s$ and $t$, we have
  \begin{equation*}
    K(s, t) = \frac{1}{2} \max(s, t) \left(\min(s, t)\right)^2 -
    \frac{1}{6} \left(\min(s, t) \right)^3. 
  \end{equation*}
\item \textbf{(scaled) Integrated Brownian Motion Plus a Linear Term}:
  In the 
  IBM model, we have $f(0) = 0$ and $f'(0) = 0$. This might be
  an unrealistic assumption to make when $f$ is completely
  unknown. In this case, a better model is:
  \begin{equation*}
    f(t) = \beta_0 + \beta_1 t + \tau \int_0^t W_s ds
  \end{equation*}
  where $\beta_0, \beta_1, \{W_s\}$ are independent with $W_s$ being
  Brownian motion and $\beta_0, \beta_1 \overset{\text{i.i.d}}{\sim} N(0,
  C)$. The kernel now becomes
  \begin{align*}
    K(s, t) &= \text{Cov} \left(\beta_0 + \beta_1 s + \tau \int_0^s W_u du,
      \beta_0 + \beta_1 t + \tau \int_0^t W_v dv \right) \\ &= C (1 +
                                                              st) +
                                                              \frac{\tau^2}{2}
                                                              \max(s,
                                                              t)
                                                              \left(\min(s,
                                                              t)\right)^2
                                                              - 
    \frac{\tau^2}{6} \left(\min(s, t) \right)^3. 
  \end{align*}
\end{enumerate}
We will see some more examples of kernels later. Next we see some
basic applications of Gaussian Processes. 


\section{Interpolation}

Suppose we are given values of $f$ at $0 \leq x_1 < \dots < x_n \leq
1$: $f(x_1), 
\dots, f(x_n)$. Let $x \in [0, 1]$ be a new point (i.e., distinct from $x_1, \dots,
x_n$). What can we say about $f(x)$?

We will solve this problem assuming a Gaussian process prior for $f$
with mean zero and covariance kernel $K(\cdot, \cdot)$. The answer
then will be given by the conditional distribution of $f(x)$ given 
$f(x_1), \dots, f(x_n)$. To compute this, we note that:
\begin{align*}
  (f(x_1), \dots, f(x_n), f(x))^T \sim N \left(0, \begin{pmatrix} (K(x_i, 
                                             x_j))_{n \times n} 
                                              & (K(x_i,
                                                            x))_{n
                                                            \times 1}
                                             \\ (K(x, x_i))_{1
                                             \times n} & K(x,
                                                         x) \end{pmatrix}
  \right). 
\end{align*}
Using the notation $K = (K(x_i, 
                                             x_j))_{n \times n}$ and
                                             $\mathbf{k} = (K(x_i,
                                                            x))_{n
                                                            \times
                                                            1}$, we
                                                          can write
                                                          the
                                                          conditional
                                                          distribution
                                                          of
                                                          $f(x)$
                                                          given $f(x_1),
                                                          \dots, 
f(x_n)$ as
\begin{align*}
  f(x) \mid f(x_1), \dots, f(x_n) \sim N\left(\mathbf{k}^TK^{-1} (f(x_1), \dots,
  f(x_n))^T, K(x, x) -
  \mathbf{k}^T K^{-1} \mathbf{k} \right).
\end{align*}
Thus the posterior mean (or mode) estimate of $f(x)$ is given by
\begin{align}\label{posmean}
  \widehat{f(x)} = \mathbf{k}^TK^{-1} (f(x_1), \dots, f(x_n))^T. 
\end{align}
For different kernels $K$, the above expression behaves differently on
$f(x_1), \dots, f(x_n)$. The simplest example is that of Brownian
Motion. 

\begin{example}[Brownian Motion]
  Suppose $f(x) = \beta_0 + \tau B(x)$ where $B(x)$ is Brownian
  Motion, $\beta_0 \sim N(0, C)$ with $C \rightarrow \infty$ and
  $\tau$ is a fixed positive constant. Then the kernel is:
  \begin{align}\label{bm_ker}
    K(x_1, x_2) = C + \tau^2 \min(x_1, x_2). 
  \end{align}
  In this case, \eqref{posmean} can be explicitly evaluated (in the
  limit $C \rightarrow \infty$) as
\begin{equation}\label{bm_interp}
\widehat{f(x)}=
\begin{cases}
f(x_1), & x \le x_1, \\[6pt]
\dfrac{x_{i+1}-x}{x_{i+1}-x_i} f(x_i)
+
\dfrac{x-x_i}{x_{i+1}-x_i} f(x_{i+1}),
& x_i \le x \le x_{i+1}, \\[10pt]
f(x_n), & x \ge x_n .
\end{cases}
\end{equation}
The details behind why \eqref{posmean} with the kernel \eqref{bm_ker}
leads to \eqref{bm_interp} will be seen later. 

The formula \eqref{bm_interp} for $\widehat{f(x)}$  performs linear
interpolation 
between the observed points, and constant extrapolation outside the
observed range.

Also note that $\widehat{f(x)}$ in \eqref{bm_interp} does not depend
on $\tau$. The 
hyperparameter $\tau$ only affects the posterior variance of $f(x)$,
and not the posterior mean. 
\end{example}

Next example is that of the Integrated Brownian Motion prior. 

\begin{example}[Integrated Brownian Motion]
  Suppose
\begin{align}\label{ibmp}
f(x)=\beta_0+\beta_1 x+\tau I(x),
\end{align}
where $I(x)$ is integrated Brownian motion and $\beta_0,\beta_1
\stackrel{i.i.d.}{\sim}N(0,C)$ with $C\to\infty$. In this case, the
posterior mean in \eqref{posmean} can be evaluated explicitly and is
the natural cubic spline interpolant of the observations. We will see
details behind this later. 
\end{example}

\section{Integration}

Suppose we are given values of $f$ at $0 \leq x_1 < \dots < x_n \leq$: $f(x_1),
f\dots, f(x_n)$. What can we say about $\int_0^1f(x)$?

Again we place a Gaussian process prior on $f$. By linearity, we can
write 
\begin{align*}
 & \E \left[ \int_0^1 f(x) dx \mid f(x_1), \dots, f(x_n) \right] \\ &=
  \int_0^1 \E \left[ f(x) \mid f(x_1), \dots, f(x_n) \right] \\ &=
                                                                  \int_0^1 \mathbf{k}^T K^{-1} (f(x_1), \dots, f(x_n))^T \\
  &= \left(\int_0^1 k(x_1, x) dx, \dots,\int_0^1 k(x_n, x) dx \right)^T
    K^{-1} (f(x_1), \dots, f(x_n))^T.  
\end{align*}

\begin{example}[Brownian Motion]
  For the Brownian motion, we already have the formula
  \eqref{bm_interp} for $\E [f(x) \mid f(x_1), \dots, f(x_n)]$. So the
  posterior mean for $\int_0^1 f(x) dx$ is obtained by simply
  integrating the right hand side of \eqref{bm_interp} from 0 to
  1. This leads to:
  \begin{align*}
    & \E \left(\int_0^1 f(x) dx \mid f(x_1), \dots, f(x_n) \right) \\
    &= x_1 f(x_1) + \sum_{i=1}^{n-1} (x_{i+1} - x_i) \frac{f(x_i) +
      f(x_{i+1})}{2} + (1 - x_n) f(x_n). 
  \end{align*}
  This the trapezoidal integration rule. Therefore, with the Brownian
  motion prior, the posterior mean estimate of $\int_0^1 f(x) dx$
  coincides with the trapezoidal rule. 
\end{example}

\begin{example}[Integrated Brownian Motion]
  With the prior \eqref{ibmp}, we get a version of the Simpson
  integration rule.  
\end{example}
For more on the relation between classical integration (quadrature)
rules and Gaussian processes, see the book
\citet{hennig2022probabilistic} or the paper
\citet{diaconis1988bayesian}.

In the  next lecture, we shall see how to use Gaussian processes for
the regression problem. 


\begin{thebibliography}{99}

\bibitem[Diaconis(1988)]{diaconis1988bayesian}
Diaconis, P. (1988)
\textit{Bayesian numerical analysis}.
Statistical Decision Theory and Related Topics IV, Vol.~1, 163--175.

%\bibitem[MacKay and Peto(1995)]{MacKay1995HierarchicalDirichlet}
%D.~J.~C. MacKay and L.~C. Peto,
%\newblock ``A hierarchical Dirichlet language model,''
%\newblock \emph{Natural Language Engineering}, vol.~1,
%no.~3, pp.~289--308, 1995.

%\bibitem[Rubin(1981)]{Rubin1981BayesianBootstrap}
%D.~B. Rubin,
%\newblock ``The Bayesian bootstrap,''
%\newblock \emph{The Annals of Statistics}, vol.~9,
%no.~1, pp.~130--134, 1981.
  
%\bibitem[Fisher(1934)]{Fisher1934}
%R.~A. Fisher,
%\newblock ``Probability, likelihood and quantity of information in the logic of
%uncertain inference,''
%\newblock \emph{Proceedings of the Royal Society of London. Series~A,
%Containing Papers of a Mathematical and Physical Character}, vol.~146,
%no.~856, pp.~1--8, 1934.

%\bibitem[Efron(2010)]{Efron}
%Efron, B. (2010)
%\textit{Large-scale inference: empirical Bayes methods for estimation,
%testing, and prediction}. Cambridge University Press, 2010.

\bibitem[Hennig et~al.(2022)]{hennig2022probabilistic}
Hennig, P., Osborne, M.~A., and Kersting, H.~P. (2022)
\textit{Probabilistic Numerics: Computation as Machine Learning}.
Cambridge University Press.

%\bibitem[Shumway and Stoffer(2017)]{ShumwayStoffer}
%Shumway, R. H. and Stoffer, D. S. (2017)
%\textit{Time Series Analysis and Its Applications: With R
%  Examples}. Springer, 4th edition. 
  
%     \bibitem[Jaynes(2003)]{Jaynes} Jaynes, E. T, (2003)
%  \textit{Probability theory: the logic of science}. Cambridge
%  University Press, 2003.
 
%  \bibitem[Sinai(1992)]{Sinai} Sinai, Y. G, (1992)
%  \textit{Probability theory: an introductory course}. Springer
%  Verlag, 1992.  

%\bibitem[{deGroot(1986)}]{DeGrootBlackwell}DeGroot, M. H. (1986)
%       A conversation with David Blackwell.
%\newblock \emph{Statistical Science}, \textbf{1}, 40--53.

%         \bibitem[Mosteller(1987)]{Mosteller} Mosteller, F, (1987)
%  \textit{Fifty challenging problems in probability with
%    solutions}. Courier Corporation, 1987.

%    \bibitem[MacKay(2003)]{MacKay} MacKay, D. J. C., (2003)
%  \textit{Information theory, inference, and learning
%  algorithms}. Cambridge University Press, 2003.

%\bibitem[{Lindley(2013)}]{Lindley}Lindley, D. V. (2013)
%   \textit{Understanding Uncertainty}. John Wiley and sons, 2013.


\end{thebibliography}







\end{document}

%%% Local Variables:
%%% mode: latex
%%% TeX-master: t
%%% End:
